\subsection{Teilmengen}
\label{subsec:fundamentals_subsets}

\begin{DefinitionBox}{Teilmenge}{fundamentals_subset}
  Eine Menge $A$ heißt \emph{Teilmenge} \EnglishTerm{subset} einer Menge $B$,
  wenn jedes Element von $A$ auch Element von $B$ ist. Man schreibt
  \[
    A\subseteq B
    \quad\Longleftrightarrow\quad
    \forall x:\; x\in A\Rightarrow x\in B.
  \]
\end{DefinitionBox}

\begin{RemarkBox}[title=\Lang{Bemerkung}{Remark}]
  \begin{minipage}[t]{0.46\linewidth}
    \textbf{Teilmenge}
    \[
      B\subseteq A
    \]
    Hier ist $B=A$ erlaubt.
  \end{minipage}
  \hfill
  \begin{minipage}[t]{0.46\linewidth}
    \textbf{Echte Teilmenge}
    \[
      B\subsetneq A
    \]
    Hier ist $B=A$ nicht erlaubt, es gilt also $B\neq A$.
  \end{minipage}
\end{RemarkBox}

Zwei Mengen sind genau dann gleich, wenn sie dieselben Elemente enthalten.
Damit gilt
\[
  A=B
  \quad\Longleftrightarrow\quad
  A\subseteq B\ \land\ B\subseteq A.
\]

\begin{ExampleBox}{Teilmengen}{fundamentals_subset_example}
  Für $A=\{1,2\}$ und $B=\{1,2,3\}$ gilt
  \[
    A\ \Cloze{1.8cm}{\subseteq B},
    \qquad
    B\ \Cloze{1.8cm}{\not\subseteq A}.
  \]
\end{ExampleBox}
